% Generated draft: requires mathematical and source review.
\documentclass[11pt]{article}
\usepackage[margin=25mm]{geometry}
\usepackage{fontspec}
\setmainfont{Latin Modern Roman}
\usepackage{amsmath,amssymb,mathtools,unicode-math}
\setmathfont{Latin Modern Math}
\usepackage{enumitem,needspace,xurl,hyperref}
\setlength{\parindent}{0pt}
\setlength{\parskip}{5pt}
\setlength{\emergencystretch}{3em}
\newcommand{\N}{\mathbb{N}}
\newcommand{\Z}{\mathbb{Z}}
\newcommand{\Q}{\mathbb{Q}}
\newcommand{\R}{\mathbb{R}}
\newcommand{\C}{\mathbb{C}}
\newcommand{\F}{\mathbb{F}}
\newcommand{\A}{\mathbb{A}}
\newcommand{\PP}{\mathbb{P}}
\newcommand{\E}{\mathbb{E}}
\newcommand{\eps}{\varepsilon}
\newcommand{\dd}{\,\mathrm{d}}
\begin{document}
\section*{UnsolvedMath ID: 8700029}
\textbf{Source problem code:} catalogue entry
\subsection{Definitions and mathematical statement}
Fix a prime $p$ and normalize $|p|_p=p^{-1}$. Let $\C_p$ be the completion of an algebraic closure of $\Q_p$, and choose an embedding of $\overline\Q$ into $\C_p$. The $p$-adic exponential is
\[
\exp_p(z)=\sum_{r=0}^{\infty}\frac{z^r}{r!},\qquad
|z|_p<p^{-1/(p-1)}.
\]
For $n\geq1$, the printed assertion says that algebraic $\beta_1,\ldots,\beta_n$ in this convergence ball have algebraically independent exponentials over $\Q$, meaning
\[
P(\exp_p(\beta_1),\ldots,\exp_p(\beta_n))\ne0
\quad\text{for every nonzero }P\in\Q[T_1,\ldots,T_n].
\]
As printed, the statement has no independence or nonzero hypothesis on the $\beta_i$. It includes $\beta_1=0$, giving $\exp_p(0)=1$, and repeated inputs, giving equal exponentials. Thus the literal assertion cannot hold.

The natural Lindemann--Weierstrass-type correction is to add $\Q$-linear independence of $\beta_1,\ldots,\beta_n$ before asking for the displayed algebraic independence. Indeed an integer linear relation among the inputs gives a multiplicative relation among their exponentials, so this is a material hypothesis, not a cosmetic convention. The correction is identified for review rather than silently substituted for the supplied statement.

Also, ``$p$-adic algebraic numbers'' must be clarified. The displayed interpretation means algebraic over $\Q$ in a chosen $p$-adic embedding. Algebraic over $\Q_p$ is a different, much larger class and cannot be used interchangeably. The full corrected claim and the supplied summary's assertions about its one-variable special case require a precise source/status check.
\subsection{Short English statement}
The intended $p$-adic Lindemann--Weierstrass question asks for algebraic independence of exponentials of small, rationally linearly independent algebraic inputs. The source omits that essential independence condition and does not specify the algebraicity base field.
\subsection{Sources}
\begin{itemize}
\item Local UnsolvedMath registry, permanent ID 8700029, exact source code catalogue entry, ``Conjecture 3.11---$p$-adic analog of Lindemann--Weierstrass's Theorem''; attributed to Michel Waldschmidt, no proposed year supplied. Registry status open, classification OPEN-TRIAGE.
\item Michel Waldschmidt, \emph{Open Diophantine Problems} (2004), Conjecture 3.11, \url{https://arxiv.org/abs/math/0312440}. Supplied source-TeX extraction accessed July 29, 2026.
\item The supplied research summary labels the algebraic-independence problem open and claims only weak progress, including an unsupported assertion that the corresponding single-exponential transcendence result is not established in full generality. No exact theorem, hypothesis range or later citation is provided; that special-case status assertion is not adopted as verified.
\item Provenance flags current primary-literature status and redistribution terms NEEDS\_REVIEW; the historical survey presents numbered open problems and conjectures.
\end{itemize}
\end{document}
